\documentclass[11pt,a4paper]{article}
\usepackage[margin=2.5cm]{geometry}
\usepackage{amsmath,amssymb,amsthm}
\usepackage{booktabs}
\usepackage{microtype}
\usepackage{xcolor}
\definecolor{cA}{HTML}{1F4E8C}
\definecolor{cB}{HTML}{B3261E}
\definecolor{cC}{HTML}{1B7A3D}
\newcommand{\seq}[1]{\texttt{\small #1}}
\newcommand{\C}{\mathcal{C}}
\newcommand{\Cp}{\mathcal{C}^{+}}
\newcommand{\Q}{\mathbb{Q}}
\newcommand{\Z}{\mathbb{Z}}
\newcommand{\R}{\mathbb{R}}
\usepackage{parskip}
\usepackage[colorlinks=true,linkcolor=blue!50!black,urlcolor=blue!50!black]{hyperref}

\theoremstyle{definition}
\newtheorem{definition}{Definition}
\theoremstyle{plain}
\newtheorem{theorem}[definition]{Theorem}
\newtheorem{lemma}[definition]{Lemma}
\newtheorem{corollary}[definition]{Corollary}
\newtheorem{proposition}[definition]{Proposition}

\title{What Morse--Hedlund means\\
\large Complexity, zero entropy, quasicrystals, and the Penrose tiling}
\author{Robin Langer}
\date{}

\begin{document}
\maketitle
\vspace{-2em}

\begin{quote}\small
This is an interpretive companion to \texttt{all-the-proofs.tex}. That document
proves things; this one asks what they mean. References of the form
\texttt{thm:mh} point into \texttt{all-the-proofs.tex}. \S\ref{sec:status}
records, claim by claim, which statements here are proved there, which are
standard results quoted from the literature, which were verified numerically
while writing this, and which are open.
\end{quote}

\section{Two theorems, and where the number $m+1$ comes from}

Write $F_\ell(m)$ for the number of distinct factors of length $m$ in a
sequence $\ell$, and recall that $\ell$ is \emph{Sturmian} when
$F_\ell(m) = m+1$ for every $m \geq 1$.

\textbf{MH-I, the dichotomy.} $F_\ell$ is bounded if and only if $\ell$ is
eventually periodic. Both halves appear in \texttt{all-the-proofs.tex}: the easy
direction as the remark that an eventually periodic sequence has bounded
complexity, and the hard direction inside the proof of
\texttt{prop:blocks}.

Unfold the dichotomy and a \emph{floor} appears. If $\ell$ is aperiodic then
$F_\ell$ is strictly increasing; since $F_\ell(1) \geq 2$, strict monotonicity
forces
\[
  F_\ell(m) \;\geq\; m+1 \qquad \text{for every } m \geq 1 .
\]
So $m+1$ is not a normalisation chosen for convenience. It is the minimum
possible complexity of an aperiodic sequence, and Sturmian sequences are exactly
those attaining it.

The mechanism behind the floor is the increment identity used in
\texttt{all-the-proofs.tex} \S7: $F_\ell(m+1) - F_\ell(m)$ is the number of
\emph{right-special} factors of length $m$, meaning words $w$ for which both
$w\seq{0}$ and $w\seq{1}$ occur. This turns the complexity function into a
derivative of the tree of factors. Aperiodicity forces at least one branch point
at every depth; Sturmian means exactly one.

\begin{quote}
Complexity is not a size. It is a branching count. A Sturmian sequence is one
whose factor tree is an infinite spine carrying a single bifurcation per level
--- the thinnest aperiodic tree there is.
\end{quote}

\textbf{MH-II, the classification.} The sequences attaining the floor are
\emph{exactly} the codings of irrational rotations. This is \texttt{thm:mh}: a
sequence is Sturmian if and only if it is $\C(\alpha,x)$ or $\Cp(\alpha,x)$ for
some irrational $\alpha \in (0,1)$ and some $x \in [0,1)$.

MH-II is the interesting one, and its interest is of a specific kind. A purely
combinatorial counting bound --- $m+1$ words of each length, not $m+2$ ---
determines a dynamical system up to two real parameters. No geometry is
assumed. No rotation, no line, no slope. They are conclusions. This is the same
shape of statement as ``a triangle-free graph with $\lfloor n^2/4 \rfloor$
edges must be complete bipartite'': saturate the inequality and lose all
freedom.

\subsection{Four readings of the same bound}

\textbf{Extremal.} Nothing in the world can sustain minimal aperiodic
complexity forever except an irrational rotation.

\textbf{Predictive.} Given the last $m$ letters, the next letter is forced for
all but one of the $m+1$ possible windows. One state of ignorance per horizon,
forever. A random binary sequence has $F(m) = 2^m$ and every window ambiguous;
Sturmian sits at the opposite end of the same axis.

\textbf{Entropic.} Topological entropy is
$h = \lim_m \log F_\ell(m) / m$, and $\log(m+1)/m \to 0$. Sturmian sequences
carry no entropy at all and are nonetheless not periodic. That combination is
what the phrase \emph{aperiodic order} names, and it is why these words are the
one-dimensional model of a quasicrystal.

\textbf{Arithmetic.} Equivalently, $F_\ell(m) = m+1$ holds precisely when $\ell$
is aperiodic and \emph{balanced}: any two factors of the same length have
$\seq{1}$-counts differing by at most $1$. Step~2 of \texttt{thm:mh} --- the
maximal $\seq{0}$-blocks take exactly two consecutive values $a-1$ and $a$ ---
is balance in local form, and balance is bounded discrepancy of the orbit
$\{k\alpha\}$. The three-distance theorem in combinatorial clothing.

\section{Zero entropy is a crude invariant}

Entropy detects only the exponential growth rate of $F_\ell$, and so collapses
an enormous amount of structure to the single value $0$.

\begin{center}
\begin{tabular}{llc}
\toprule
$F_\ell(m)$ & class & $h$ \\
\midrule
bounded & eventually periodic & $0$ \\
\textcolor{cA}{$\boldsymbol{m+1}$} & \textcolor{cA}{\textbf{Sturmian}} & $0$ \\
$km + c$ & quasi-Sturmian, Pisot substitutive & $0$ \\
$\sim 10m/3$ & Thue--Morse & $0$ \\
$\sim m^2$ & polygonal billiards, two-dimensional Sturmian & $0$ \\
$2^{m/\log m}$ & subexponential zero-entropy systems & $0$ \\
$2^{hm}$ & positive entropy & $h > 0$ \\
\bottomrule
\end{tabular}
\end{center}

Every row but the last is ``zero entropy''. So the observation that Sturmian
sequences have zero entropy is true and nearly vacuous. The content is that they
sit at the bottom of that column and that \emph{the bottom is a single point}:
$m+1$ is attained, and attaining it is a complete classification. Nothing else
in the table is rigid in this way. ``Linear complexity'' is a large and messy
class; ``complexity exactly $m+1$'' is two real parameters.

\begin{quote}
Entropy is $\lim \log F / m$; complexity is $F$. Taking the logarithm and the
limit is exactly the step that destroys rigidity. Morse--Hedlund is interesting
because it works with the unlogged object, which is also why the proof of
\texttt{thm:mh} runs on the increment identity rather than on any entropy
argument. Entropy cannot see a difference of $1$.
\end{quote}

\section{What replaces entropy: the spectrum}

Since the entropy vanishes, the classifying invariants are spectral. Let $X_\alpha$
be the Sturmian subshift, the closure of the shift orbit of $\ell$.

\begin{itemize}
\item $X_\alpha$ is \emph{uniquely ergodic}: one invariant measure, so factor
  frequencies exist and are forced. The frequency of $\seq{1}$ is $\alpha$,
  which is \texttt{cor:slope}.
\item Measure-theoretically $(X_\alpha, \text{shift})$ is isomorphic to the
  circle under rotation by $\alpha$. Its Kolmogorov--Sinai entropy is $0$ and
  its spectrum is \emph{pure discrete}, with eigenvalues
  $\{e^{2\pi i n\alpha} : n \in \Z\}$. It is a Kronecker system: almost
  periodic in Bohr's sense.
\item Topologically it is \emph{not} the rotation. The coding map
  $\text{circle} \to X_\alpha$ is a bijection except over the countably many
  points whose orbit meets the partition endpoint $0$, where the fibre has
  exactly two elements --- and those two elements are $\C(\alpha,x)$ and
  $\Cp(\alpha,x)$.
\end{itemize}

That last bullet is \texttt{thm:counterex} stated dynamically. The sequence
$\seq{1}c$ of that theorem is one of the two preimages of a single point of the
circle. In the standard terminology $X_\alpha$ is \emph{almost automorphic}: an
almost one-to-one extension of its maximal equicontinuous factor.

\begin{quote}
The disjunction in \texttt{thm:mh} is therefore not a blemish. It is the name of
the only way in which $X_\alpha$ fails to be a rotation. The circle is the
geometry; $X_\alpha$ is the combinatorics; they differ exactly on one countable
orbit, where the floor and ceiling conventions split. A measure-zero
disagreement, invisible to entropy and invisible to $F_\ell$, and the precise
reason the naive form of the theorem is false.
\end{quote}

\section{Quasicrystals: the same construction in two dimensions}

Read a Sturmian word geometrically. Replace $\seq{1}$ by a long interval $L$ and
$\seq{0}$ by a short one $S$. For the Fibonacci word, $|L|/|S| = \varphi$ and the
result is the \emph{Fibonacci chain}, the standard one-dimensional quasicrystal
of condensed-matter physics. By \texttt{thm:fib} it is the fixed point of
$\sigma(\seq{1}) = \seq{10}$, $\sigma(\seq{0}) = \seq{1}$, which is $\tau_1$ in
the notation of \texttt{all-the-proofs.tex} \S5.

Where does it come from? From the cutting-sequence picture, which is the
degenerate case of \emph{cut-and-project}: take a lattice, a subspace of
irrational slope, a strip, and project the lattice points inside the strip.

\begin{center}\small
\begin{tabular}{lll}
\toprule
 & Fibonacci chain & Penrose tiling \\
\midrule
lattice & $\Z^2$ & $\Z^5$ (really $A_4$: $\Z^5$ modulo its diagonal) \\
physical space $E_\parallel$ & line of slope $\alpha$ & $2$-plane invariant under $5$-fold rotation \\
internal space $E_\perp$ & complementary line & complementary $2$-plane \\
window & interval, the projected unit square & pentagons, the projected unit $5$-cube \\
window decomposes into & $2$ subintervals $\leftrightarrow$ letters & $4$ pentagons $\leftrightarrow$ vertex classes \\
parameter & intercept $x \in [0,1)$ & offsets $\gamma_1,\dots,\gamma_5$, $\sum \gamma_j \in \Z$ \\
aperiodic because & $\alpha$ irrational & $E_\parallel$ totally irrational \\
inflation factor & $\varphi$ & $\varphi$ \\
\bottomrule
\end{tabular}
\end{center}

The same machine, run from dimension $2$ to $1$ and from dimension $5$ to $2$.
The physical quasicrystals --- Shechtman's 1982 Al--Mn alloy with fivefold
diffraction symmetry, Nobel Prize 2011 --- are the planar and spatial cases.
``Fibonacci is to Penrose as $\Z^2$ is to $\Z^5$'' is not an analogy; it is the
definition in both cases.

\section{The Penrose tiling contains Sturmian words}

Decorate each Penrose rhomb in the standard way and the segments join, across
the whole tiling, into five families of infinite straight lines at angles
$2\pi k/5$: the \emph{Ammann bars}. Within one family the lines are not equally
spaced. The gaps take exactly two values $L$ and $S$ with $L/S = \varphi$, and
the sequence of gaps read along the perpendicular direction is the Fibonacci
word.

A Penrose tiling is therefore, to a first approximation, five Fibonacci words
laid at $72^\circ$ to one another, subject to a consistency condition. Penrose's
matching rules are equivalent to the requirement that the bars continue straight
across every edge. The aperiodicity is enforced in five one-dimensional
channels, each an instance of the object in \texttt{all-the-proofs.tex}.

This yields an aperiodicity proof straight from MH-I.

\begin{proposition}\label{prop:aper}
A Penrose tiling has no translation period. The same argument covers the
Ammann--Beenker and dodecagonal tilings.
\end{proposition}

\begin{proof}
Suppose $v$ were a period. Then each Ammann family's gap sequence would be
periodic, hence of bounded factor complexity by the easy half of MH-I. But that
gap sequence is the Sturmian word of slope $1/\varphi$, whose complexity is
$m+1$ and so unbounded.

The Ammann-bar fact is quoted. Here is the same conclusion using only the grid
formalism of \S6.2, which is proved. A translation by $v$ carries the family-$j$
lines $\{p : p \cdot e_j + \gamma_j \in \Z\}$ to themselves only if
$v \cdot e_j \in \Z$, and it must do so for every $j$, since the tiling
determines its grid. Fix two families $m$ and $j$: then
$(v \cdot e_j)/(v \cdot e_m)$ is the ratio $\alpha$ of \S6.2, which is
irrational --- it is $\pm\varphi^{\pm 1}$, $\pm\sqrt 2^{\pm 1}$ or
$\pm\sqrt 3^{\pm 1}$ in the three cases. A ratio of two integers cannot be
irrational unless both vanish, so $v \cdot e_j = 0$ for all $j$, and the $e_j$
span $\R^2$, so $v = 0$.
\end{proof}

So the irrationality of the measured slopes \emph{is} the aperiodicity: the same
number that makes the index word Sturmian is the obstruction to a period. A
search over all $v$ with $v \cdot e_0, v \cdot e_1 \in \{-8,\dots,8\}$ found no
nonzero candidate for $n = 4, 5, 6$.

One distinction worth preserving: de Bruijn's \emph{pentagrid} is a different
five-family construction, in which the lines \emph{are} equally spaced within
each family and the tiling is the dual of the grid. There the irrationality
enters through the incommensurability of the five directions rather than through
the spacings. Both pictures are correct; the Ammann bars are the one in which
Sturmian words appear verbatim.

\subsection{Why the golden ratio --- twice, for the same reason}

Fivefold symmetry puts the tiling's coefficients in the cyclotomic field
$\Q(\zeta_5)$, whose maximal real subfield is
\[
  \Q(\zeta_5) \cap \R \;=\; \Q(\sqrt 5) \;=\; \Q(\varphi) .
\]
So $\varphi$ is forced the moment fivefold symmetry is demanded.

Meanwhile, in \texttt{all-the-proofs.tex} \S8, $\varphi$ arises for a
completely different-looking reason: $a_1 = 1$ and $\beta = \alpha$, that is,
$\alpha = 1/\varphi$ is the fixed point of the Gauss map
$\alpha \mapsto \{1/\alpha\}$. That is precisely why derivation returns the same
sequence and the infinite renormalisation collapses to a single substitution.

These are the same number reached from crystallography and from continued
fractions, and the coincidence is not one. A Sturmian word is substitutive ---
self-similar under a single inflation, rather than under an infinite sequence of
distinct ones --- exactly when its slope is a quadratic irrational, i.e.\ has
eventually periodic continued fraction. A quasicrystal needs a self-similar
inflation in order to have matching rules at all; quadratic irrationality is
the price; and the lowest-degree option carrying a non-crystallographic rotation
is $\varphi$. Hence the Fibonacci chain in physics, hence pentagons in Penrose,
and hence the fact that \S8 of \texttt{all-the-proofs.tex} is the one place in
the paper where the Gauss orbit is a fixed point.

\subsection{Inflation, Pisot, diffraction}

\begin{itemize}
\item \textbf{Inflation.} Penrose tilings inflate by $\varphi$: each rhomb
  subdivides into smaller rhombs and the result is the original rescaled.
  \texttt{thm:fib}'s $\sigma(c) = c$ and $c = \lim_j \sigma^j(\seq{1})$ is the
  one-dimensional instance --- the tiling is its own inflation, grown from a
  seed.
\item \textbf{Pisot.} Penrose's substitution matrix has leading eigenvalue
  $\varphi^2$ with its conjugate inside the unit circle; $\tau_1$'s matrix
  $\left(\begin{smallmatrix}1&1\\1&0\end{smallmatrix}\right)$ has $\varphi$
  and $-1/\varphi$. Both Pisot, so both have pure point diffraction.
\item \textbf{Diffraction.} Fibonacci: Bragg peaks on the rank-$2$ module
  $\Z + \Z\varphi$. Penrose: Bragg peaks on a rank-$4$ module, with tenfold
  symmetric intensities --- the Shechtman pattern. Projecting the Penrose
  Fourier module along one Ammann direction returns $\Z + \Z\varphi$.
\item \textbf{Frequencies.} The thick-to-thin rhomb ratio is $\varphi$, forced
  by the eigenvector of the inflation matrix. The one-dimensional statement is
  \texttt{cor:slope} evaluated at $\alpha = 1/\varphi$.
\end{itemize}

\subsection{Conway worms are the Penrose form of $\C$ versus $\Cp$}

This parallel is exact, and it is the one I find most convincing.

In \texttt{thm:counterex} the coding map is injective except over the countable
orbit of $0$, where the fibre has two elements: floor and ceiling, because the
point sits on the partition boundary and a side must be chosen.

In the Penrose family the parameter is the offset vector $\gamma$. For
\emph{generic} $\gamma$ the projection data determine the tiling uniquely. For
\emph{singular} $\gamma$ --- where lattice points land on the window's boundary
--- the construction is ambiguous, and the ambiguity is resolved along a line:
one obtains distinct tilings with identical finite patches, differing by flips
along an infinite strip. Those strips are \emph{Conway worms}; the maximally
symmetric singular configuration is the \emph{decapod}, whose tiling is Conway's
cartwheel.

In both cases: the hull is an almost one-to-one extension of a torus; the
non-injectivity lives over a measure-zero set of boundary-hitting parameters;
and no finitely observable datum --- the language, the patch set, $F_\ell$, the
diffraction pattern --- distinguishes members of a singular fibre. The
disjunction in \texttt{thm:mh} and the existence of Conway worms are one
phenomenon in two dimensions' worth of clothing.

This also explains how there can be uncountably many Penrose tilings while
``the'' Penrose tiling remains a meaningful phrase: all of them are locally
indistinguishable, every finite patch of one occurring in every other. The
Sturmian statement is that all $\C(\alpha,x)$ with $\alpha$ fixed share one
language, so $F_\ell$ depends on $\alpha$ only through its irrationality and on
$x$ not at all. Complexity sees the hull, never the point in it.

\section{Other quadratic irrationals}

Fibonacci is one Sturmian word among uncountably many. Which of the others
give Penrose-like objects? The question splits in two, with very different
answers in dimension one and dimension two.

\subsection{Dimension one: the metallic means, out of $\tau_a$}

Abelianise the substitution $\tau_a(\seq{1}) = \seq{1}\seq{0}^a$,
$\tau_a(\seq{0}) = \seq{1}\seq{0}^{a-1}$ of
\texttt{all-the-proofs.tex} \S5, in the basis (count of $\seq{1}$s, count of
$\seq{0}$s):
\[
  M_a \;=\; \begin{pmatrix} 1 & 1 \\ a & a-1 \end{pmatrix},
  \qquad \operatorname{tr} M_a = a, \qquad \det M_a = -1 ,
\]
so the inflation factor satisfies
\[
  \lambda^2 - a\lambda - 1 = 0, \qquad
  \lambda \;=\; \frac{a + \sqrt{a^2+4}}{2} .
\]
These are the \emph{metallic means}: golden at $a=1$, silver $1+\sqrt 2$ at
$a=2$, bronze $(3+\sqrt{13})/2$ at $a=3$, and onward. The silver-mean chain,
$\tau_2$ with $\seq{1} \mapsto \seq{100}$ and $\seq{0} \mapsto \seq{10}$, is a
standard object in the quasicrystal literature alongside Fibonacci.

\begin{quote}
$\det M_a = -1$ for every $a$, and any product of such matrices has determinant
$\pm 1$. So the inflation factor of a self-similar Sturmian chain is always an
algebraic \emph{unit} of a real quadratic field, with conjugate $\mp 1/\lambda$
of modulus below $1$ --- hence always Pisot, automatically. ``Quadratic slope
gives a good quasicrystal'' is not luck; it is $\det = \pm 1$, which is itself
only the statement that $\tau_a$ is invertible on the abelianisation because the
block decomposition \texttt{lem:blocks} is a bijection.
\end{quote}

Eventually periodic continued fractions supply the rest. A periodic block
$[0;\overline{a_1,\dots,a_p}]$ yields the composite substitution
$\tau_{a_1} \circ \cdots \circ \tau_{a_p}$, whose matrix product has the
quadratic unit as leading eigenvalue. For instance
$M_1 M_2 = \left(\begin{smallmatrix}3&2\\1&1\end{smallmatrix}\right)$ has trace
$4$ and determinant $1$, giving $\lambda^2 - 4\lambda + 1 = 0$ and
$\lambda = 2+\sqrt 3$.

The fixed point itself is pinned, and its intercept is \emph{not} the slope
except in the Fibonacci case. That is worth recording, because it is exactly
what makes \texttt{thm:fib} look simpler than the general statement.

\begin{proposition}\label{prop:metallic}
Fix $a \geq 1$ and let $\alpha = \bigl(\sqrt{a^2+4} - a\bigr)/2 = 1/\lambda$,
the irrational with constant partial quotients $[0;\overline{a}]$. Then the
unique Sturmian word fixed by $\tau_a$ is
\[
  \C\!\left(\alpha, \tfrac{1}{1+\alpha}\right) .
\]
For $a = 1$ one has $\alpha^2 + \alpha = 1$, hence $1/(1+\alpha) = \alpha$, and
the statement reduces to \texttt{thm:fib}.
\end{proposition}

\begin{proof}
The slope condition is $\{1/\alpha\} = \alpha$, i.e.\ $1/\alpha = a + \alpha$,
i.e.\ $\alpha^2 + a\alpha - 1 = 0$, which is the stated $\alpha$; this is the
fixed point of the Gauss map with partial quotient $a$, so $a_1 = a$ and
$\beta = \alpha$ as required by \texttt{cor:gauss}.

Since $\tau_a(\seq{1})$ and $\tau_a(\seq{0})$ both begin with $\seq{1}$, the
fixed point begins with $\seq{1}$, so $k_0 = 0$ and $x_0 = x$. The coordinate
computation in the proof of \texttt{thm:fib} reads, in general,
$u_0 = \bigl(x - (1-\alpha)\bigr)/\alpha$, and \texttt{thm:derive} gives the
derived intercept as $\{-u_0\}$. Demanding $\{-u_0\} = x$ means
$-(x - 1 + \alpha)/\alpha = x + m$ for some $m \in \Z$, that is
\[
  x \;=\; \frac{1 - (m+1)\alpha}{1+\alpha} .
\]
The constraint $k_0 = 0$ is $x \geq 1 - \alpha$. For $m = 0$ we get
$(1-\alpha)/(1+\alpha) < 1-\alpha$, excluded; for $m = -2$ we get $x = 1$,
outside $[0,1)$; $m \leq -3$ and $m \geq 1$ leave $[0,1)$ likewise. For
$m = -1$, $x = 1/(1+\alpha)$ and $1/(1+\alpha) \geq 1-\alpha$ holds because
$(1-\alpha)(1+\alpha) = 1 - \alpha^2 \leq 1$. So $m = -1$ is the only
admissible choice.
\end{proof}

Verified numerically for $a = 1, \dots, 7$ by comparing the first $8000$ letters
of the $\tau_a$-fixed point against $\C(\alpha, 1/(1+\alpha))$, agreement exact
in every case. The composite $\tau_1 \tau_2$ behaves the same way with its own
intercept: its fixed point is
$\C\!\left(\sqrt 3 - 1,\ (\sqrt 3 - 1)/2\right)$, agreeing on the first
$50\,000$ letters.

The clean general statement of which slopes admit a substitution is the theory
of \emph{Sturm numbers}: $\alpha \in (0,1)$ is the slope of a
substitution-invariant Sturmian word if and only if $\alpha$ is a quadratic
irrational whose algebraic conjugate $\alpha^*$ does not lie in $(0,1)$
(Crisp--Moran--Pollington--Shiue; Allauzen). Quadratic is necessary, and the
point on the orbit must then be chosen correctly ---
Proposition~\ref{prop:metallic} is that choice, made explicit, for constant
partial quotients.

\subsection{Dimension two: exactly three fields, because $\phi(n) = 4$}

Now impose $n$-fold rotational symmetry on a planar tiling. Its coefficients
must lie in the cyclotomic ring $\Z[\zeta_n]$, of rank $\phi(n)$ over $\Z$, so
the cut-and-project must run from $\Z^{\phi(n)}$ down to $2$ dimensions. The
inflation factor is real and lives in the maximal real subfield
$\Q(\zeta_n) \cap \R$, of degree $\phi(n)/2$.

For the inflation factor to be \emph{quadratic} we therefore need
$\phi(n)/2 = 2$, that is $\phi(n) = 4$, that is
\[
  n \in \{5, 8, 10, 12\} .
\]
That list is complete, and it is exactly the roster of classical planar
quasicrystals.

\begin{center}
\begin{tabular}{llllp{4.6cm}}
\toprule
$n$ & field & $\lambda$ & tiling & matching Sturmian slope \\
\midrule
$5, 10$ & $\Q(\sqrt 5)$ & $\varphi$ & \textbf{Penrose} & $[0;\overline{1}] = 1/\varphi$, from $\tau_1$ \\
$8$ & $\Q(\sqrt 2)$ & $1+\sqrt 2$ & \textbf{Ammann--Beenker} & $[0;\overline{2}] = \sqrt 2 - 1$, from $\tau_2$ \\
$12$ & $\Q(\sqrt 3)$ & $2+\sqrt 3$ & \textbf{Stampfli, shield} & $[0;\overline{1,2}] = \sqrt 3 - 1$, from $\tau_1\tau_2$ \\
\midrule
$7,9,14,18$ & cubic & $\approx 2.247$ & Danzer's sevenfold & none: no quadratic slope \\
\bottomrule
\end{tabular}
\end{center}

So the answer in its strongest form: \emph{the only other quadratic irrationals
giving Penrose-like planar tilings are $\sqrt 2$ and $\sqrt 3$.} Each is
Penrose-like in every structural sense --- since $\phi(n) = 4$ in all three
cases, each is a codimension-two projection of a rank-four module, $\Z^4 \to \R^2$,
onto an $n$-fold invariant plane, with an inflation symmetry, matching rules and
pure point diffraction with $n$-fold Bragg peaks --- and each carries its own
Sturmian word along its bar families, with the two spacings in the ratio
supplied by that field's unit. Note that the three rows of the table are three
rows of the $\tau_a$ tower of \texttt{all-the-proofs.tex}: $\tau_1$, $\tau_2$,
and $\tau_1 \tau_2$.

\subsubsection*{The slopes, measured rather than quoted}

The right-hand column of that table was quoted. It can be measured. Write the
$n$ grid families as $\{p : p \cdot e_j + \gamma_j \in \Z\}$ with
$e_j = e^{2\pi i j/n}$ for odd $n$ and $e^{\pi i j/n}$ for even $n$, and recall
de Bruijn's index
\[
  K_j(p) \;=\; \lceil\, p \cdot e_j + \gamma_j \,\rceil .
\]
This is the cumulative counting function of \texttt{lem:count}, so its
increments are a coding. Walk a line of family $k$, in direction
$d = e_k^\perp$, sampling at its crossings with family $m$ --- which are equally
spaced, and so supply the index of a coding --- and count with family $j$. Then
\[
  K_j(\text{$i{+}1$st sample}) - K_j(\text{$i$th sample})
  \;=\; \Cp(\alpha, c)_i + \lfloor \alpha \rfloor ,
  \qquad
  \alpha \;=\; \frac{d \cdot e_j}{d \cdot e_m} ,
\]
the $\lfloor \alpha \rfloor$ because
$\lceil i\alpha + x \rceil = \lceil i\{\alpha\} + x\rceil + i\lfloor\alpha\rfloor$.
So every grid line of every one of these tilings carries a coding, at a slope
fixed by the geometry alone.

Computing those slopes gives

\begin{center}
\begin{tabular}{llll}
\toprule
$n$ & tiling & $\{\alpha\}$ measured & continued fraction \\
\midrule
$5$ & Penrose & $0.618034 = 1/\varphi$, $0.381966 = 1/\varphi^2$ & $[0;\overline{1}]$, $[0;2,\overline{1}]$ \\
$4$ & Ammann--Beenker & $0.414214 = \sqrt 2 - 1$ & $[0;\overline{2}]$ \\
$6$ & dodecagonal & $0.732051 = \sqrt 3 - 1$ & $[0;\overline{1,2}]$ \\
\bottomrule
\end{tabular}
\end{center}

which are exactly the slopes of $\tau_1$, $\tau_2$ and $\tau_1\tau_2$ from
\S\ref{sec:status}'s metallic-mean family. The agreement is not built in: the
left side comes from pentagonal, octagonal and dodecagonal geometry, the right
side from the Gauss map. In every case tested the increments agreed with
$\Cp(\alpha,c)$ letter for letter, and the measured complexity was
$F(m) = m+1$ for $m = 1,\dots,8$ --- so the tilings' own grid indices are
Sturmian, measured rather than assumed.

Every quasicrystal so far found in a real material is quadratic: icosahedral
(the original Al--Mn) is $\Q(\sqrt 5)$, decagonal is $\Q(\sqrt 5)$, octagonal
is $\Q(\sqrt 2)$, dodecagonal is $\Q(\sqrt 3)$. Danzer's sevenfold tilings exist
mathematically, and their cubic inflation factor
$1 + 2\cos(2\pi/7) \approx 2.2470$ is even Pisot --- its conjugates
$1+2\cos(4\pi/7) \approx 0.5550$ and $1+2\cos(6\pi/7) \approx -0.8019$ both have
modulus below $1$. Yet no sevenfold alloy has ever been observed, in four
decades of looking. The plausible reason is the next subsection.

\subsection{$\Q(\sqrt 2)$ is a twin of $\Q(\sqrt 5)$; $\Q(\sqrt 3)$ is not}

The three fields of \S6.2 are not on an equal footing, and the asymmetry is
forced by arithmetic rather than by geometry.

\begin{center}
\begin{tabular}{llllll}
\toprule
field & $n$ & $\lambda$ & minimal polynomial & $N(\lambda)$ & substitution \\
\midrule
$\Q(\sqrt 5)$ & $5, 10$ & $\varphi$      & $x^2 - x - 1$  & $-1$ & $\tau_1$ \\
$\Q(\sqrt 2)$ & $8$     & $1 + \sqrt 2$  & $x^2 - 2x - 1$ & $-1$ & $\tau_2$ \\
$\Q(\sqrt 3)$ & $12$    & $2 + \sqrt 3$  & $x^2 - 4x + 1$ & $+1$ & $\tau_1\tau_2$ \\
\bottomrule
\end{tabular}
\end{center}

$\Q(\sqrt 2)$ is the exact twin of $\Q(\sqrt 5)$: its unit is the metallic mean
at $a = 2$, so $\lambda^2 = 2\lambda + 1$, a single $\tau_a$ suffices, and the
slope has constant partial quotients. $\Q(\sqrt 3)$ is not, and cannot be.

\begin{proposition}\label{prop:norm}
Let $\alpha$ have purely periodic continued fraction of period $p$ and let
$\lambda$ be the leading eigenvalue of $M_{a_1} \cdots M_{a_p}$. Then
\[
  N(\lambda) \;=\; \det\bigl(M_{a_1} \cdots M_{a_p}\bigr) \;=\; (-1)^p .
\]
Consequently, if $\Q(\sqrt d)$ contains no unit of norm $-1$, then every
Sturmian slope generating $\Q(\sqrt d)$ has \emph{even} continued-fraction
period. In particular no single $\tau_a$ has inflation factor in $\Q(\sqrt 3)$,
and the dodecagonal case needs a period at least $2$.
\end{proposition}

\begin{proof}
For a $2 \times 2$ integer matrix the determinant is the product of the
eigenvalues, which for $M_{a_1}\cdots M_{a_p}$ are $\lambda$ and its algebraic
conjugate; so $N(\lambda) = \det$. Each $\det M_{a} = -1$ by \S6.1, whence
$\det = (-1)^p$. If $p$ were odd then $\lambda$ would be a unit of norm $-1$.

For $d = 3$ no such unit exists: $x^2 - 3y^2 = -1$ requires
$x^2 \equiv -1 \equiv 2 \pmod 3$, and the squares modulo $3$ are $0$ and $1$.
The direct form of the same obstruction: a single $\tau_a$ generates
$\Q(\sqrt{a^2+4})$, and $a^2 + 4 = 3k^2$ would give $a^2 \equiv 2 \pmod 3$,
again impossible.
\end{proof}

By contrast $x^2 - 2y^2 = -1$ has the solution $(1,1)$ and $x^2 - 5y^2 = -1$ has
$(2,1)$, which is why those two fields are reached at period $1$, by $\tau_2$ and
$\tau_1$ respectively.

Proposition~\ref{prop:norm} is visible in the measured slopes. Every slope read
off the Penrose and Ammann--Beenker tilings has continued fraction of period
$1$ --- $[0;\overline{1}]$, $[0;2,\overline{1}]$ for the first,
$[0;\overline{2}]$, $[0;1,\overline{2}]$, $[0;3,\overline{2}]$ for the second
--- while every slope read off the dodecagonal tiling has period $2$:
$[0;\overline{1,2}]$, $[0;1,\overline{1,2}]$, $[0;2,\overline{2,1}]$. The
twelvefold quasicrystal is one storey taller in the tower of \S\S4--5 than the
other two, and Pell's equation is the reason.

\subsection{Quadratic is the boundary, and the boundary is the Gauss map}

The deep statement concerns \emph{local rules}. A tiling has local matching
rules when a finite set of permitted neighbour configurations forces membership
in the hull. That is what makes a quasicrystal growable, atom by atom, rather
than merely definable.

\begin{itemize}
\item \textbf{One dimension.} A Sturmian word is substitution-invariant, hence
  locally rule-generated, if and only if its slope is a Sturm number --- in
  particular quadratic.
\item \textbf{Two dimensions.} A planar cut-and-project tiling admits local
  rules only if the projection subspace is defined over a quadratic field
  (Le; related results of Levitov and Socolar). Beyond quadratic, no finite rule
  set suffices.
\end{itemize}

In the language of \texttt{all-the-proofs.tex} \S\S4--5 the reason is
immediate. Derivation sends $\alpha \mapsto \{1/\alpha\}$ and strips a block of
width $a_i$. Matching rules amount to needing only finitely many distinct
$\tau_a$, which is to say the Gauss orbit of $\alpha$ must be eventually
periodic, which by Lagrange is to say $\alpha$ is quadratic. For
$\alpha = e - 2$, or any Liouville number, the partial quotients never cycle,
each storey of the tower demands a fresh substitution with a fresh block width,
and no finite local rule can see arbitrarily long blocks.

\begin{quote}
\texttt{cor:gauss} --- the block widths encountered are the partial quotients of
$\alpha$ --- is therefore the quasicrystal-quality classifier. Finite Gauss
orbit, $\alpha$ rational: periodic, a crystal. Eventually periodic, $\alpha$
quadratic: self-similar, matching rules, a genuine quasicrystal, and in the
plane one of exactly four rotational symmetries. Aperiodic Gauss orbit,
everything else: still a perfectly good quasiperiodic tiling with sharp Bragg
peaks, but no inflation and no local rules. Definable, not growable. The
derivation operator of \S4 is the thing that decides which.
\end{quote}

A related and sharper quality measure: a Sturmian word is \emph{linearly
recurrent} --- every factor of length $m$ recurs within $Cm$ --- if and only if
$\alpha$ is of bounded type, $\sup_i a_i < \infty$. Lagarias and Pleasants argue
that linear recurrence is the right definition of an \emph{ideal} quasicrystal.
Liouville slopes, with unbounded partial quotients, produce arbitrarily long
stretches that look locally periodic.

\section{What quadratic buys, and what it does not}

A distinction easy to blur. Pure point diffraction does \emph{not} require
quadratic, or even Pisot. Every regular model set --- cut-and-project with any
totally irrational subspace and a well-behaved window --- diffracts purely
pointwise (Hof, Schlottmann). So $\C(\alpha,x)$ has sharp Bragg peaks on
$\Z + \Z\alpha$ for \emph{every} irrational $\alpha$.

The Pisot condition matters for the other family: systems defined by a
substitution rather than by a projection. There pure point diffraction can fail.
\textbf{Thue--Morse} is the standard example --- complexity about $10m/3$, hence
zero entropy and low complexity by any reasonable standard, yet singular
continuous diffraction, no Bragg peaks at all. So zero entropy does not buy
quasicrystalline order; the Pisot property of the inflation factor does, and for
quadratic slopes it comes free by $\det M_a = -1$.

Quadratic irrationality therefore buys self-similarity and local rules. It does
not buy the diffraction, which was already there.

The resulting stratification of ``Penrose-like'':
\begin{enumerate}
\item \textbf{Quasiperiodic with pure point diffraction.} Every irrational
  $\alpha$. Uncountably many, one per slope.
\item \textbf{Self-similar, substitutive, locally rule-generated.} Quadratic
  $\alpha$ only. Countably many; $\tau_a$ and its products generate them all.
\item \textbf{Planar, with a crystallographically forbidden rotational symmetry
  and quadratic inflation.} Exactly $n = 5, 8, 10, 12$. Penrose,
  Ammann--Beenker, dodecagonal. Three fields: $\Q(\sqrt 5)$, $\Q(\sqrt 2)$,
  $\Q(\sqrt 3)$.
\end{enumerate}

\section{The physics, briefly}

Put an electron on the chain. The discrete Schr\"odinger operator
\[
  (H\psi)_k \;=\; \psi_{k+1} + \psi_{k-1} + V \ell_k \psi_k ,
  \qquad \ell \text{ Sturmian},
\]
is the \emph{Fibonacci Hamiltonian}. The three complexity regimes give three
different physics.

\begin{center}
\begin{tabular}{lll}
\toprule
potential & spectrum & transport \\
\midrule
periodic ($F_\ell$ bounded) & bands, absolutely continuous & ballistic, $|x| \sim t$ \\
\textcolor{cA}{Sturmian ($F_\ell = m+1$)} & \textcolor{cA}{zero-measure Cantor set,} & \textcolor{cA}{anomalous, $|x| \sim t^\gamma$,} \\
 & \textcolor{cA}{purely singular continuous} & \textcolor{cA}{$0 < \gamma < 1$} \\
i.i.d.\ random (positive entropy) & pure point, dense & Anderson localisation, none \\
\bottomrule
\end{tabular}
\end{center}

For every $V \neq 0$ the spectrum is a Cantor set of zero Lebesgue measure
(S\"ut\H{o}; Bellissard--Iochum--Scoppola--Testard; the fractal-dimension theory
is Damanik--Gorodetski--Yessen). This is the mathematical reason quasicrystalline
alloys conduct so startlingly badly for metals: their electronic states are
neither extended nor localised.

And the control parameter is the continued fraction of $\alpha$. The
transfer-matrix renormalisation is a \emph{trace map} indexed by the partial
quotients --- the derivation operator of \S4 run on matrices instead of words.
The Hausdorff dimension of the spectrum and the transport exponent $\gamma$
depend on $\alpha$ through the $a_i$, and the same Diophantine dichotomy
governs the almost Mathieu operator, the Hofstadter butterfly, and the Ten
Martini problem (Avila--Jitomirskaya).

\section{Where the analogy stops}

Everything above transfers from dimension one to dimension two. The
\emph{theorem} does not. There is no known two-dimensional Morse--Hedlund.

For a planar configuration let $P(m,n)$ count distinct $m \times n$ patches.
Penrose has $P(m,n) \sim C m^2 n^2$: polynomial, zero entropy, but nothing like
a rigid extremal value. The conjectured analogue of MH-I is \textbf{Nivat's
conjecture} (1997):

\begin{quote}
If $P(m,n) \leq mn$ for some $m, n \geq 1$, then the configuration is periodic.
\end{quote}

Open in general after nearly thirty years. Partial results: Kari--Szabados for
algebraic and low-complexity cases, Cyr--Kra for $P(m,n) \leq mn/2$ and related
bounds, with various strengthenings by Szabados. And there is no known
two-dimensional analogue of MH-II at all --- no classification saying that the
configurations attaining the minimum are exactly the projections of
such-and-such a lattice. The Penrose tiling is not extremal for any known
complexity functional. It is extremal for \emph{symmetry} instead: the
lowest-degree field admitting a non-crystallographic rotation.

Which is the honest answer to why MH-II deserves a careful proof, and a
formalisation. One dimension is the only place where the quasicrystal story has
a complete converse, and \texttt{thm:mh} together with the
$\C$-or-$\Cp$ repair is that converse stated correctly.

\section{Status of the claims}\label{sec:status}

\textbf{Proved in \texttt{all-the-proofs.tex}.} MH-I in both directions; MH-II
(\texttt{thm:mh}); the increment identity; the counterexample
\texttt{thm:counterex} and the $\C$-or-$\Cp$ repair; \texttt{cor:slope},
\texttt{cor:gauss}, \texttt{lem:blocks}, \texttt{thm:derive},
\texttt{thm:fib}.

\textbf{Proved here.} Proposition~\ref{prop:aper}, in two ways: from MH-I plus
the quoted Ammann-bar fact, and --- independently of anything quoted --- from
the grid formalism of \S6.2, covering all three symmetries.
Proposition~\ref{prop:metallic} (the $\tau_a$ fixed point is
$\C(\alpha, 1/(1+\alpha))$), from \texttt{thm:derive}.
Proposition~\ref{prop:norm} ($N(\lambda) = (-1)^p$, so $\Q(\sqrt 3)$ needs even
period), from $\det M_a = -1$ and a congruence modulo $3$.

\textbf{Verified numerically while writing this.} $\phi(n) = 4$ exactly for
$n \in \{5,8,10,12\}$. The metallic-mean eigenvalues and their conjugates.
Proposition~\ref{prop:metallic} for $a = 1,\dots,7$ on $8000$-letter prefixes.
The $\tau_1\tau_2$ fixed point as $\C(\sqrt 3 - 1, (\sqrt 3 - 1)/2)$ on
$50\,000$-letter prefixes. That $1+2\cos(2\pi/7)$ is Pisot.

The measured-slope table of \S6.2, by an independent implementation of the
generalised dual method (\texttt{penrose-codings.html}, with its kernel written
as $\lceil k\alpha + x\rceil$ and its increments): that the construction yields
an exact tiling for $n = 4,5,6$ (every sample point covered exactly once, no
gaps, no overlaps); that the index increments equal $\Cp(\alpha,c)$ letter for
letter; that $F(m) = m+1$ for $m \leq 8$; that for $n = 5$ de Bruijn's index
sums stay bounded while the thick-to-thin ratio approaches $\varphi$.

That $\Q(\sqrt 3)$ is never $\Q(\sqrt{a^2+4})$, checked for $a \leq 60$; that
$x^2 - 3y^2 = -1$ has no solution with $y \leq 2000$ while $d = 2, 5$ have small
ones; the continued-fraction periods of all eight measured slopes; and that no
nonzero $v$ with $v \cdot e_0, v \cdot e_1 \in \{-8,\dots,8\}$ satisfies
$v \cdot e_j \in \Z$ for all $j$, for $n = 4, 5, 6$.

\textbf{Quoted from the literature, not verified here.} The Ammann-bar
description of Penrose and the corresponding bar families of
Ammann--Beenker; unique ergodicity and the almost automorphic structure of
$X_\alpha$; Conway worms and the singular Penrose tilings; the cut-and-project
descriptions and the $\phi(n)$ rank argument; pure point diffraction of regular
model sets; the local-rules theorems; the Sturm-number characterisation; the
Fibonacci-Hamiltonian spectral results.

\textbf{Open.} Nivat's conjecture. Any two-dimensional analogue of MH-II. The
absence of sevenfold quasicrystalline alloys has no proof, only the
local-rules heuristic of \S6.3.

\textbf{Attempted and withdrawn.} The identification of a single phason flip
with a single $\C$-versus-$\Cp$ choice at one index. Numerically, perturbing
$\gamma$ through the maximally singular value $\gamma = 0$ changes roughly half
the tiles in a patch, spread over the whole patch (principal-axis ratio $1.25$),
not confined to a strip. That is consistent with $\gamma = 0$ being the
tenfold-degenerate decapod configuration rather than a single worm, so the
experiment does not refute \S5.3 --- but it does not support the sharper
one-index claim either, and \S5.3 should be read as the structural analogy it
states and nothing more.

\section{Sources}

\begin{itemize}\itemsep2pt
\item M.~Baake and U.~Grimm, \emph{Aperiodic Order}, vol.~1, CUP 2013. \S7.1
  Fibonacci, \S7.3 Penrose, \S7.8 Ammann--Beenker, written as parallel sections
  precisely to expose the dictionary of \S4 above; chs.~4 and 9 for diffraction.
\item N.~G.~de Bruijn, ``Algebraic theory of Penrose's non-periodic tilings of
  the plane'', I and II, \emph{Indag.\ Math.} 43 (1981). The pentagrid, and the
  Ammann bars.
\item M.~Senechal, \emph{Quasicrystals and Geometry}, CUP 1995. The singular
  tilings, worms, and the cartwheel.
\item N.~Pytheas Fogg, \emph{Substitutions in Dynamics, Arithmetics and
  Combinatorics}, Springer LNM 1794, 2002. Ch.~6 is the Sturmian/continued
  fraction dictionary in the form closest to \texttt{all-the-proofs.tex}
  \S\S4--8.
\item T.~T.~Q.~Le, ``Local rules for quasiperiodic tilings'', in
  \emph{The Mathematics of Long-Range Aperiodic Order}, Kluwer 1997. The
  quadratic-field theorem.
\item D.~Crisp, W.~Moran, A.~Pollington, P.~Shiue, ``Substitution invariant
  cutting sequences'', \emph{J.\ Th\'eor.\ Nombres Bordeaux} 5 (1993);
  C.~Allauzen, ``Une caract\'erisation simple des nombres de Sturm'',
  \emph{J.\ Th\'eor.\ Nombres Bordeaux} 10 (1998).
\item D.~Damanik, J.~Fillman, A.~Gorodetski, surveys on the Fibonacci
  Hamiltonian; A.~Avila and S.~Jitomirskaya, ``The Ten Martini Problem'',
  \emph{Ann.\ of Math.} 170 (2009).
\item M.~Nivat, ICALP 1997 invited talk (the conjecture); J.~Kari and
  M.~Szabados, V.~Cyr and B.~Kra, for the partial results.
\end{itemize}

\end{document}
